% 图 4-54（题 4-18）：按原书图中的轨迹、日期点及比例尺同比例描绘。
% 原图比例尺 0--30 的长度约 105 像素；轨道两端约相距 169 像素。
% 本图用于原题的读图估算，日期点沿原图保留，不作额外轨道拟合。
\begin{tikzpicture}[>=Stealth,thick,x=0.055cm,y=-0.055cm,font=\small]
  \draw (0,0) .. controls (-7,-7) and (2,-5) .. (12,2)
    --(132,79) .. controls (135,81) and (137,83) .. (138,84)
    .. controls (146,92) and (147,98) .. (139,92)
    --(109,80)--(79,61)--(59,45)--(36,30)--(15,17)--(3,6)
    .. controls (-1,3) and (-3,0) .. (0,0);
  \fill (1,2) circle (1.2pt) node[above right=3pt] {木星中心};
  \foreach \x/\y/\lx/\ly/\label in {
    3/6/-17/24/7月15日,
    15/17/-2/41/7月1日,
    36/30/8/58/6月15日,
    59/45/28/75/5月15日,
    79/61/47/92/3月30日,
    109/80/65/110/1994年1月27日,
    139/92/161/102/7月1日,
    138/84/162/74/1993年3月25日} {
      \fill (\x,\y) rectangle ++(2,2);
      \draw[thin] (\x,\y)--(\lx,\ly);
      \node[fill=white,inner sep=1pt] at (\lx,\ly) {\label};
  }
  \draw[thin] (-25,145)--(90,145);
  \foreach \x/\label in {-20/0,15/10,50/20,85/30} {
    \draw[thin] (\x,142)--(\x,148);
    \node[below=3pt] at (\x,148) {\label};
  }
  \node[below=16pt] at (32.5,148) {距离 / $10^6\,\mathrm{km}$};
\end{tikzpicture}
